% cap2_studii.tex
\chapter{Studii anterioare}
\label{ch:studii}

Prezentul capitol trece în revistă stadiul actual al domeniilor relevante pentru această
lucrare: platformele de automatizare rezidențială, protocoalele de comunicație \gls{iot},
paradigmele de calcul la margine (edge/fog) și metodele de detecție a anomaliilor comportamentale.
Pe baza acestei analize, sunt identificate lacunele din literatura de specialitate pe care
sistemul propus urmărește să le acopere.

% ======================================================================
\section{Platforme smart home comerciale și open-source}
\label{sec:studii-platforme}
% ======================================================================

Piața automatizărilor rezidențiale a cunoscut o expansiune rapidă în ultimul deceniu, atât
în segmentul comercial cât și în cel open-source. Stojkoska și Trivodaliev \cite{stojkoska2017smarthome}
oferă o trecere în revistă cuprinzătoare a provocărilor și soluțiilor din domeniu, evidențiind
că interoperabilitatea, securitatea și confidențialitatea datelor rămân problemele centrale
nerezolvate pe deplin de nicio platformă.

\subsection{Platforme comerciale}

Printre soluțiile comerciale cu cea mai largă adopție se numără Amazon Alexa/Echo,
Google Home și Apple HomeKit. Toate trei se bazează pe arhitecturi hibride cloud–dispozitiv:
comenzile utilizatorului sunt transmise serverelor producătorului, procesate în cloud și
returnate dispozitivului. Această abordare oferă avantaje clare — actualizări automate,
recunoaștere vocală de înaltă calitate, ecosisteme bogate de aplicații — dar implică o
serie de dezavantaje fundamentale. Alam et al.\ \cite{alam2012smarthome} subliniază că
dependența de conectivitate la internet face ca astfel de sisteme să devină nefuncționale
în absența conexiunii sau la întreruperea serviciului cloud al producătorului.

Aspectul confidențialității este tratat în detaliu de Zheng et al.\ \cite{zheng2018smarthomeprivacy},
care arată că datele colectate de dispozitivele smart home comerciale — inclusiv obiceiuri
zilnice, prezența la domiciliu sau conversații ambientale — sunt transferate și stocate pe
servere terțe, deseori în jurisdicții cu standarde de protecție a datelor diferite. O
analiză de securitate sistematică realizată de Alrawi et al.\ \cite{alrawi2019sokhome} pe
un eșantion larg de dispozitive IoT de consum a identificat vulnerabilități la nivelul
firmware-ului, al protocoalelor de comunicație și al \gls{api}-urilor cloud, concluzionând
că nivelul de securitate al ecosistemelor comerciale rămâne nesatisfăcător în raport cu
amenințările actuale.

\subsection{Platforme open-source}

La polul opus se situează platformele open-source, dintre care Home Assistant și openHAB
sunt cele mai reprezentative. Home Assistant, scris în Python și distribuit sub licență
Apache 2.0, rulează local pe hardware modest (inclusiv Raspberry Pi) și oferă o interfață
web completă, suport pentru sute de integrări și un motor de automatizare bazat pe
evenimente. openHAB, scris în Java, adoptă o filosofie similară de independență față de
cloud, cu accent pe extensibilitate prin plugin-uri.

Avantajul esențial al acestor platforme este că întreaga logică de control și toate datele
rămân în rețeaua locală a utilizatorului. Totuși, curba de învățare este semnificativă,
iar suportul pentru \gls{ml} nativ (mai ales pentru detecția de anomalii) este limitat sau
realizat prin integrări externe.

Sistemul propus în această lucrare împărtășește filosofia local-first a platformelor
open-source, dar se distinge prin integrarea unui motor de \gls{ml} propriu — Isolation
Forest — rulat direct pe gateway-ul Raspberry Pi, fără a depinde de nicio bibliotecă
externă de automatizare generică.

% ======================================================================
\section{Protocoale de comunicație IoT}
\label{sec:studii-protocoale}
% ======================================================================

Alegerea protocolului de comunicație are implicații directe asupra consumului de resurse,
a latenței, a scalabilității și a fiabilității unui sistem \gls{iot}. Cele mai utilizate
protocoale în domeniu sunt \gls{mqtt}, CoAP și HTTP. Naik \cite{naik2017mqttcomparison}
realizează o comparație detaliată a acestora, evidențiind că nu există o soluție universală
optimă, alegerea depinzând de constrângerile concrete ale aplicației.

\subsection{Comparație MQTT / CoAP / HTTP}

Tabelul~\ref{tab:protocoale} sintetizează principalele caracteristici ale celor trei
protocoale din perspectiva aplicabilității \gls{iot}.

\begin{table}[H]
\centering
\caption{Comparație MQTT, CoAP și HTTP pentru aplicații IoT}
\label{tab:protocoale}
\begin{tabular}{@{}p{3.2cm} p{3.8cm} p{3.8cm} p{3.8cm}@{}}
\toprule
\textbf{Caracteristică} & \textbf{MQTT} & \textbf{CoAP} & \textbf{HTTP/1.1} \\
\midrule
Transport            & TCP           & UDP           & TCP \\
\addlinespace
Model comunicație    & Publish/Subscribe (broker central) & Request/Response (P2P sau proxy) & Request/Response (client–server) \\
\addlinespace
Overhead antet       & Minim (2–5 B antet fix) & Redus (4 B antet fix) & Mare (text, zeci–sute de octeți) \\
\addlinespace
QoS                  & 3 niveluri (0, 1, 2) & 2 tipuri (confirmabil / neconfirmabil) & Implicit (TCP garantează livrarea) \\
\addlinespace
Last Will / Retained & Da (suport nativ) & Nu (necesită extensii) & Nu \\
\addlinespace
Potrivire IoT        & Excelentă (multe-la-mulți, noduri cu resurse limitate) & Bună (rețele cu pierderi, IPv6/6LoWPAN) & Moderată (overhead ridicat, nepotrivit pentru microcontrolere) \\
\addlinespace
Securitate           & TLS/SSL opțional & DTLS opțional & TLS/SSL standard \\
\addlinespace
Standard             & OASIS \cite{oasis2019mqtt} & IETF RFC 7252 \cite{rfc7252} & IETF RFC 7230 \\
\bottomrule
\end{tabular}
\end{table}

\subsection{Justificarea alegerii MQTT}

\gls{mqtt} a fost conceput explicit pentru conexiuni cu lățime de bandă mică și dispozitive
cu resurse limitate \cite{oasis2019mqtt}. Modelul publish/subscribe decuplează producătorul
de date (nodul ESP32) de consumatori (backend NestJS, motor ML), permitând mai multor
abonați să recepționeze simultan aceleași date fără ca nodul să cunoască numărul sau
identitatea acestora.

Caracteristica \emph{Last Will and Testament} este esențială pentru monitorizarea stării
nodurilor: la conectare, fiecare nod înregistrează un mesaj de „ultimă voință" pe topicul
\texttt{smarthome/\{nodeId\}/status}; dacă conexiunea se întrerupe neașteptat, brokerul
publică automat acest mesaj, informând restul sistemului. Mesajele \emph{retained} asigură
că un abonat nou primește imediat ultima valoare a unui topic, fără a aștepta o nouă
publicare.

Mosquitto, implementarea de referință pentru medii embedded \cite{light2017mosquitto},
rulează eficient pe Raspberry Pi 3B+ și suportă autentificare prin nume de utilizator și
parolă, acoperind cerințele de securitate minimale ale rețelei locale. Combinația MQTT +
Mosquitto este larg utilizată în proiecte IoT de cercetare \cite{naik2017mqttcomparison},
ceea ce facilitează reproductibilitatea experimentelor.

CoAP ar fi o alternativă pentru rețele cu pierderi mari sau IPv6/6LoWPAN, scenarii
irelevante pentru o rețea Wi-Fi locală. HTTP, deși simplu de implementat, ar impune un
overhead de zeci de ori mai mare față de \gls{mqtt} pentru aceeași sarcină utilă de date
senzoriale, epuizând rapid memoria stack-ului TCP de pe microcontroler.

% ======================================================================
\section{Edge computing și procesarea locală}
\label{sec:studii-edge}
% ======================================================================

Paradigma \emph{edge computing} plasează capacitatea de calcul cât mai aproape de sursa
de date, reducând latența, consumul de lățime de bandă și dependența de conectivitate la
internet. Shi et al.\ \cite{shi2016edge} definesc edge computing ca extensie a modelului
cloud la marginea rețelei și argumentează că este esențial pentru aplicații cu cerințe
stricte de latență și confidențialitate.

În contextul \gls{iot}, Premsankar et al.\ \cite{premsankar2018edgeiot} demonstrează că
procesarea locală reduce semnificativ volumul de date transmis spre cloud și permite
reacții în timp real imposibil de obținut printr-o arhitectură cu cloud central. Lin et
al.\ \cite{lin2017edge} extind analiza la scenarii de calcul mobil, arătând că un nod
edge (cum ar fi un Raspberry Pi) poate prelua sarcini de inferență \gls{ml} de complexitate
medie cu latențe sub o secundă.

Sistemul proiectat în această lucrare adoptă integral filozofia edge computing: Raspberry
Pi 3B+ acționează ca nod edge care stochează toate datele local (PostgreSQL + TimescaleDB),
execută inferența Isolation Forest în timp real și publică alerte fără a implica niciun
serviciu extern. Această arhitectură garantează funcționarea completă a sistemului chiar
și în absența totală a internetului, răspunzând direct preocupărilor de autonomie și
confidențialitate identificate în literatura de specialitate.

% ======================================================================
\section{Detecția de anomalii comportamentale}
\label{sec:studii-anomalii}
% ======================================================================

Detecția de anomalii reprezintă procesul de identificare a observațiilor care se abat
semnificativ de la comportamentul așteptat al unui sistem. Chandola et al.\
\cite{chandola2009anomaly} oferă o taxonomie exhaustivă a metodelor existente, grupate
în funcție de tipul de supraveghere și de tehnica statistică utilizată.

\subsection{Taxonomia metodelor de detecție}

\paragraph{Metode statistice.} Presupun că datele normale urmează o distribuție
probabilistică parametrică (Gaussiană, Poisson etc.). Punctele cu probabilitate mică
sunt considerate anomalii. Avantajul principal este interpretabilitatea matematică,
dar metodele parametrice eșuează când distribuția reală a datelor este necunoscută
sau multimodală — situație frecventă în datele senzoriale din locuințe.

\paragraph{Metode bazate pe distanță și densitate.} \emph{Local Outlier Factor} (LOF),
propus de Breunig et al.\ \cite{breunig2000lof}, compară densitatea locală a unui punct
cu cea a vecinilor săi. LOF este eficient pentru clustere de densități diferite, dar
costul computațional $\mathcal{O}(n^2)$ îl face nepotrivit pentru inferență în timp real
pe dispozitive edge cu resurse limitate.

\paragraph{One-Class SVM.} Scholkopf et al.\ \cite{scholkopf2001ocsvm} propun o variantă
de \gls{svm} antrenată exclusiv pe date normale, care delimitează o frontieră de decizie
în spațiul de caracteristici. Performanța depinde critic de alegerea kernelului și a
hiperparametrilor, iar costul de inferență crește cu numărul de vectori suport.

\paragraph{Metode de izolare.} Liu, Ting și Zhou \cite{liu2008isolation} introduc
\emph{Isolation Forest}, un algoritm nesupervizat care izolează punctele anomale prin
partiționări aleatoare succesive ale spațiului de caracteristici. Ideea centrală este
că un punct anomal — rar și îndepărtat — va fi separat de restul printr-un număr mic
de partiții, în timp ce un punct normal, aflat în zona densă, necesită multe partiții
pentru a fi izolat. Scorul de anomalie este invers proporțional cu lungimea medie a
căii de izolare. Liu et al.\ \cite{liu2012isolationtkdd} confirmă eficiența algoritmului
pe seturi de date de mari dimensiuni, cu complexitate $\mathcal{O}(n \log n)$ la
antrenare și $\mathcal{O}(\log n)$ la inferență.

\paragraph{Metode de învățare profundă.} Chalapathy și Chawla \cite{chalapathy2019deepanomaly}
trec în revistă utilizarea rețelelor neuronale profunde (autoencoder, LSTM, GAN) pentru
detecția de anomalii. Aceste metode ating performanțe superioare pe seturi mari de date
cu structuri complexe, dar necesită resurse computaționale semnificative (GPU, memorie
RAM de ordinul gigaocteților) și date etichetate sau volume mari de antrenare, cerințe
incompatibile cu un nod edge de tip Raspberry Pi 3B+.

Cook et al.\ \cite{cook2020anomalyiot} și Fahim și Sillitti \cite{fahim2019anomaly}
examinează specific tehnicile de detecție a anomaliilor în fluxuri \gls{iot}, subliniind
necesitatea unor metode incrementale sau cu footprint redus. Concluzia lor converge cu
alegerea Isolation Forest pentru sisteme cu constrângeri de resurse.

\subsection{Principiul izolării — ilustrare conceptuală}

\figref{fig:ml-isolation} ilustrează intuitiv diferența dintre comportamentul unui punct
normal și cel al unui punct anomal în algoritmul Isolation Forest. Un punct anomal,
aflat în afara norului dens de observații normale, este separat prin numai câteva
partiționări aleatoare. Prin contrast, un punct normal, înconjurat de alți vecini apropiați,
necesită un număr mult mai mare de partiții înainte de a fi izolat. Această asimetrie
stă la baza scorului de anomalie: cu cât calea de izolare este mai scurtă, cu atât
probabilitatea ca observația să fie o anomalie este mai mare.

\begin{figure}[H]
  \centering
  \includegraphics[width=\doublefigure]{\pics/fig-ml-isolation}
  \caption{Principiul izolării: un punct anomal este separat prin puține partiționări
    aleatoare, în timp ce un punct normal (din zona densă) necesită multe partiții.}
  \label{fig:ml-isolation}
\end{figure}

\subsection{Justificarea alegerii Isolation Forest pentru sistem}

Alegerea Isolation Forest ca algoritm de bază al motorului \gls{ml} din această lucrare
este motivată de trei proprietăți esențiale, deosebit de relevante pentru mediile edge:

\begin{enumerate}
  \item \textbf{Nesupervizat.} Nu necesită date etichetate cu anomalii cunoscute — o
    cerință imposibil de îndeplinit în practică pentru o locuință nouă, unde anomaliile
    reale sunt rare și nedefinite a priori.

  \item \textbf{Eficiență computațională.} Complexitatea $\mathcal{O}(\log n)$ la
    inferență \cite{liu2012isolationtkdd} permite scorarea în timp real a fiecărui set
    de citiri senzoriale pe un Raspberry Pi 3B+ (ARMv8, 1 GB RAM), fără a afecta
    restul serviciilor (broker MQTT, baza de date, \gls{api}).

  \item \textbf{Consum redus de memorie.} Modelul serializat (\texttt{isolation\_forest.pkl})
    ocupă câțiva zeci de kiloocteți, spre deosebire de rețelele neuronale profunde care
    pot necesita sute de megaocteți. Implementarea din scikit-learn \cite{pedregosa2011scikit}
    este matură, testată exhaustiv și bine documentată.
\end{enumerate}

Pe lângă aceste argumente tehnice, Isolation Forest s-a dovedit robust pe date
multidimensionale cu distribuții arbitrare \cite{liu2008isolation}, ceea ce este important
deoarece vectorul de caracteristici al sistemului combină valori continue (temperatură,
umiditate, gaz), valori normalizate logaritmic (luminozitate) și variabile binare/ciclice
(oră din zi, zi din săptămână, indicator noapte/weekend).

% ======================================================================
\section{Poziționarea lucrării față de stadiul actual}
\label{sec:studii-pozitionare}
% ======================================================================

Trecerea în revistă de mai sus permite identificarea clară a nișei pe care o ocupă
sistemul propus față de soluțiile existente.

\textbf{Față de platformele comerciale} (Alexa, Google Home, HomeKit), sistemul propus
elimină dependența de cloud și garantează că niciun datum personal nu părăsește rețeaua
locală. Funcționarea completă în regim offline și costul de infrastructură zero (nu există
abonamente sau servicii de hosting) constituie avantaje practice concrete pentru utilizatorul
final, în spiritul principiilor local-first formalizate de Kleppmann et al.\
\cite{kleppmann2019localfirst}.

\textbf{Față de platformele open-source generice} (Home Assistant, openHAB), sistemul
propus integrează nativ un motor de detecție a anomaliilor comportamentale antrenat pe
datele specifice locuinței, nu un sistem de reguli fix sau o integrare externă opțională.
Această abordare permite identificarea comportamentelor anormale subtile — de exemplu, un
model de temperatură neobișnuit la o anumită oră din zi — imposibil de captat prin simple
praguri hardcodate.

\textbf{Față de lucrările de cercetare pe detecția de anomalii IoT} \cite{cook2020anomalyiot,
fahim2019anomaly}, contribuția principală a prezentei lucrări constă în implementarea
completă end-to-end a unui pipeline funcțional pe hardware real de cost redus: de la
colectarea datelor senzoriale printr-un nod ESP32, prin stocarea în TimescaleDB, la
antrenarea zilnică automată a unui model Isolation Forest și inferența în timp real pe
Raspberry Pi 3B+. Lucrările citate tratează frecvent problema detecției de anomalii
în izolare, fără a aborda provocările de integrare sistem sau constrângerile edge specifice.

\textbf{Noutatea esențială} a lucrării este combinarea a trei caracteristici rareori
întâlnite simultan în literatura de specialitate: (1)~arhitectură complet \emph{local-first},
fără nicio dependență de servicii externe; (2)~detecție de anomalii comportamentale
prin \gls{ml} nesupervizat (\emph{Isolation Forest}), antrenat pe profilul real al
locuinței; (3)~implementare pe hardware de consum redus și cost accesibil (ESP32 + Raspberry
Pi), cu o interfață mobilă React Native pentru utilizatorul final.

Capitolul următor (\chapref{ch:metodologie}) detaliază arhitectura sistemului și
deciziile de proiectare care decurg din analiza prezentată mai sus.
